\section{Tax information before tax calculation}
\label{sec:shared-tax-information}

The mixed book has recorded profit, but it still lacks a tax assessment.
The earlier return-only reference does contain a tax-close method. Whether
that method can be reused is a substantive question: the new book contains
lease rights and a modification gain that the old book did not represent.
The unchanged method refuses the mixed state because its closed state shape
does not admit those additional accounts and lease records. The two original
tax examples still reproduce their complete results. Thus the refusal identifies
an incompatible extension, rather than a regression in the original method.

\subsection{A new income category changes the tax starting point}

The original pretax function sums the original income accounts, excluding
current and deferred tax. If that function is applied diagnostically to the
mixed balances, it returns a loss of 15. The complete recorded report instead
includes the lease-modification gain of 20. Under the signed-account convention,
the additional \texttt{other\_gains} balance is $-20$, so
\[
  P_{\mathrm{old\ accounts}}=-15,
  \qquad
  P_{\mathrm{mixed}}=-15-(-20)=5.
\]
The difference is an omitted income category, not a rounding error. Replacing
the state validator while leaving the original income list unchanged would
therefore calculate the wrong pretax starting point for the stated mixed-book
target. It would not preserve the meaning of the old calculation.

Even the corrected starting point does not determine tax. The return example
has a signed reconciliation adjustment of $-24$. Mechanical reuse of that
adjustment and the original illustrative rate would give
\[
  \frac{-15-24}{4}=-\frac{39}{4},
  \qquad
  \frac{5-24}{4}=-\frac{19}{4}.
\]
These expressions diagnose the effect of the omitted gain; neither is an
assessed tax or a cash refund. The lease deduction timing, the tax treatment
of the modification, the remaining right-of-use tax bases and the treatment
of a loss have not been supplied. The original finite method refuses its
negative-current-tax route. Copying the rate or reversing every deferred-tax
balance would not supply the missing facts.

Dates create a separate incompatibility. The original close on 31 December
stores 1 January as an exclusive tax endpoint. The mixed-book invariant
currently requires that endpoint to be no later than the state date. Those
two meanings cannot be combined by changing only the endpoint field. Moreover,
the mixed example's execution window runs from 1 January 2024 to 1 January 2026:
731 days. That window contains the selected lease and return events, but does
not establish a single annual tax interval or complete fiscal activity.

\subsection{Keep missing declarations and rejected premises visible}

A separate information assessor now accepts the unchanged recorded-report
request, an explicitly stated tax interval and identified information
declarations. The interval distinguishes its start, its last included calendar
day and its exclusive endpoint. The first finite implementation requires the
start to equal the recorded opening and every requested event to lie inside
the interval. The last recorded event may precede the asserted coverage date;
that assertion remains a declaration, not evidence that all activity is known.

For this mixed example, the assessor identifies 61 required declarations. They
cover jurisdiction and regime, the interval basis and coverage findings, rate,
loss and cash-settlement treatment, deferred recognition and recovery, every
pretax account's treatment and signed reconciliation adjustment, balance-sheet
tax bases, and the tax bases and timing of each lease right. Zero balances
remain in the account population. Both lease rights A and B remain relevant:
the native state retains fully terminated B with zero asset and liability,
and its earlier modification still requires a stated treatment.

Each supplied value carries an assumption or source-reference description.
Each unknown carries its reason. A source reference is retained as supplied;
the assessor does not authenticate it. With no information supplied, all
61 declarations remain unresolved. Removing any one from the hypothetical
complete control again leaves an unresolved requirement. A false coverage or
recognition finding remains false and is reported as unsupported. Extra fields,
wrong scalar types, forged request identities and incompatible interval
boundaries are invalid.

The complete native report accompanies the assessment. Its stopped event
prefix, component memory, gross cash movements and unresolved annual coverage
are preserved. When a requested payment fails, supplying every information
declaration does not erase that failure: completion of the requested sequence
remains unresolved. A first-year interval can be stated separately when its
requested events fit that interval, without relabelling the two-year example.

\subsection{Supplied information is not an executable tax rule}

The positive result is deliberately called \texttt{information\_supplied}.
It means the finite declaration set is present and has the required types.
The hypothetical controls exercise that representation; their supplied zero
bases and adjustments are not inferred from the financial statements. A text
field describing a deduction can still fail to specify a valid calculation.
Likewise, individually typed amounts can be inconsistent when aggregated.

Accordingly, every result keeps calculation readiness, source and accounting
admission, tax execution and annual-close readiness false. No tax amount is
posted. The next calculation design must state how account adjustments,
right-level tax bases, deferred recognition and settlements relate, and must
check their joint consistency. This preserves useful progress: missing facts
and competing declarations can be represented and examined now, while the
remaining numerical and accounting questions remain explicit.

\subsection{Reconcile a declared current-tax calculation}

The next finite calculation makes one further assumption explicit: a single
nonnegative rate applies to nonnegative taxable profit before credits or any
other regime-specific adjustments. This is a hypothetical arithmetic model,
named \texttt{uniform\_rate\_no\_credits\_v1}. It does not select a country's
tax rules. The author must select that model affirmatively; a false selection
is unsupported and an unknown selection remains unresolved.

Each declared account adjustment is now supported by identified allocation
items. An item taken from the completed return history must retain its original
event identity and signed amount, and must appear exactly once. Additional
items have explicit assumption references. The items assigned to an income
account must sum to that account's declared adjustment. Their category assignment
is supplied by the author; the arithmetic check does not establish its legal
appropriateness. In particular, the native return adjustment is not added again
after it has already been included in those account totals.

There are also four cross-field basis checks. At both opening and closing,
the declared asset bases of all lease rights must sum to the aggregate
right-of-use tax basis, and the declared liability bases must sum to the
aggregate lease-liability basis. The zero-valued right B participates in those
checks. A mismatch rejects the conditional calculation while preserving the
complete recorded financial report and the inconsistent declarations for
inspection. Neither carrying amount nor termination status supplies a tax basis.

\begin{samepage}
Let $P$ be recorded pretax income over all seven income accounts, $a_i$ the
identified signed adjustment items and $r$ the explicitly supplied rate. The
conditional taxable-profit expression and gross current tax are
\[
 T=P+\sum_i a_i,
 \qquad
 C=rT\quad\text{when }T\geq0.
\]
\end{samepage}
Here $C$ is gross conditional current tax under the selected arithmetic rule.
It is not an incremental expense, accrued payable or cash payment. The
calculation does not subtract opening accruals, apply credits, recognize
deferred tax or change the recorded closing state.

For the mixed example, $P=5$ and the retained return adjustment is $-24$.
Without any additional declared item, $T=5-24=-19$. The result is
\texttt{loss\_route\_required}; no negative tax amount or refund is computed.
A text declaration requesting a refund cannot change that result because it
does not supply an executable loss rule. Two hypothetical controls illustrate
the nonnegative branch. An expressly supplied additional adjustment of 19 gives
$T=0$ and $C=0$. An additional adjustment of 40 gives $T=21$ and, at the
illustrative rate $1/4$, $C=21/4$. Those extra adjustments are test assumptions,
not inferred company facts or recommended tax treatments.

The comparison retains the entire financial report, including cash of 361 and
recorded profit of 5, in every accepted arithmetic result. Inconsistent allocation
or basis declarations retain that same report alongside their rejection. Tests
also change native adjustment amounts, duplicate or omit identities, alter the
model finding and tamper with computed results. These checks establish the
specified arithmetic and preservation properties. A loss-utilization rule,
tax settlement, deferred recognition and annual close remain separate work.

\subsection{Carry a loss quantity without inventing a tax asset}

The negative expression of $-19$ needs a rule before it can affect a later
tax calculation. The finite extension supplies a conditional loss ledger.
Each lot has an identity, an origin reference, an availability date, an explicit
expiry declaration, original and remaining quantities, cumulative usage and
expiry, and an eligibility finding. An unknown finding remains unresolved;
a false finding is unsupported. A generated loss therefore becomes a quantity
in this separate ledger only under an affirmative, traceable assumption.
It is not a cash refund or a recognized deferred-tax asset.

The selected illustrative policy evaluates all declared usage on the last
included calendar day of an interval. An exclusive expiry on that day makes
the lot unavailable for usage. Expiry at the following endpoint permits usage
on the preceding day, but the remainder expires before the next opening.
No expiry is itself an explicit assumption. These conventions specify this
finite model; another jurisdiction or an intraperiod method may require a
different rule. Neither FIFO nor a legal utilization percentage is assumed.
The author supplies the chosen lot allocations and an interval usage cap.

For each lot, original loss equals remaining, used and expired quantities.
Across a step, with $L_0$ and $L_1$ denoting opening and closing remaining
quantities, $G$ generation, $U$ usage and $E$ expiry,
\[
 L_0+G=U+E+L_1.
\]
All quantities on this line are nonnegative exact rationals. A generated loss
becomes available at the exclusive endpoint and cannot be used in its own
generation interval. A usage cannot exceed its lot's remaining quantity,
the declared interval cap or positive taxable profit. Zero and exhausted lots
remain identified, preventing the same identity from being introduced again.
Consecutive ledger entries must agree on both interval endpoints and their
opening and closing quantities.

The checked negative example generates 19 under explicit eligibility and
no-expiry premises. A separate hypothetical 2026 calculation supplies taxable
profit of 21. Explicit use of 19 leaves taxable profit of 2 and gross conditional
tax of $1/2$ at rate $1/4$. This connects the loss ledger across periods; it does
not establish financial-book carry. The 2026 financial opening is independently
declared for the test and is expressly not the carried closing state of the
2024--2026 example. Both complete conditional calculations remain attached.

Partial use makes the expiry rule consequential. Using 7 of 19 leaves 12 when
the expiry is later than the next endpoint. If the exclusive expiry equals that
endpoint, 7 may be used on the preceding day and the remaining 12 expires.
If expiry instead equals the usage day, the usage is rejected; with no usage,
all 19 expires. These outcomes follow the selected policy, rather than a
silent decision to use the oldest lot or to preserve an expired balance.

The operation proposes a complete next loss state and per-lot movements.
Applying it recomputes the proposal and compares the complete result, including
the original financial calculation. Changed amounts, origins, eligibility,
scope, cursors or even Boolean-to-integer substitutions are rejected. This is
a pure in-memory application: it creates no durable commit, tax journal,
deferred asset or annual authorization. Connecting the conditional tax amount
and loss ledger to an actual financial tax balance still requires separately
specified recognition, accrual and settlement methods.

\subsection{Accrue current tax before choosing its payment}

A conditional calculation supplies an amount, but recognition and payment
change different financial balances. The next finite profile makes those
changes explicit. It begins with the complete recorded return-and-lease
calculation, recomputes the conditional loss result and requires an affirmative
recognition premise. An unknown premise remains unresolved; a false premise
does not produce a financial state. The profile applies only when the native
period has no current-tax expense already recorded and its tax-close cursor
still equals the assessment start. It therefore does not guess how an opening
tax accrual or an earlier assessment should be allocated.

The successor adds a separate income-tax payable to the original 25 accounts.
That liability is not trade payable. If the recomputed gross conditional tax
after loss usage is $C$, the accrual debits current-tax expense and credits
income-tax payable by $C$. Cash does not change. A zero amount produces a zero
journal. Existing deferred-tax assets, deferred-tax liabilities and deferred-tax
expense remain exactly as recorded; the loss balance does not create a new
deferred asset.

Payment requires its own identity, date, positive amount and affirmative
instruction. It debits income-tax payable and credits cash, subject to both the
outstanding liability and available cash. The author supplies the amount;
neither full payment nor a refund is inferred. A shortage of cash blocks that
payment while leaving a valid accrual available. This illustrative profile
classifies its specified tax payments as operating cash movements.

\begin{samepage}
Let $S$ be cumulative payment against the newly accrued amount $C$, with
$0\leq S\leq C$. Relative to the original recorded statements, the changes are
\[
 \Delta\text{profit}=-C,\qquad
 \Delta\text{tax payable}=C-S,\qquad
 \Delta\text{cash}=-S.
\]
\end{samepage}
Assets decrease by $S$, liabilities increase by $C-S$ and equity decreases by
$C$, preserving the balance-sheet identity. In the indirect operating cash
bridge, the lower profit and the increase in tax payable contribute
$-C+(C-S)=-S$. This equals the separately recorded cash payment. Paying the
liability does not charge the same expense again.

For the separate hypothetical 2026 calculation, loss usage leaves conditional
tax of $1/2$. Accrual records expense and payable of $1/2$ while cash stays at
500. An explicit payment of $1/4$ leaves cash of $1999/4$ and payable of $1/4$;
a second payment of $1/4$ clears the payable and leaves cash of $999/2$.
The recorded deferred-tax asset of 10 and liability of 4 are unchanged.
The original mixed example instead generates the loss quantity 19 and a zero
tax journal. These are consequences of the declared synthetic inputs, not
assessments of either a company or a jurisdiction.

The financial date and assessment coverage have separate meanings. Accrual is
dated on the last included day; coverage ends at the following exclusive
endpoint. The successor records that endpoint separately and leaves the native
tax-close cursor unchanged. Payments cannot precede the last tax operation or
fall outside the original financial period. An end-of-period accrual therefore
cannot silently accept a payment in the next year: financial-period carry must
first be specified.

The original native schema refuses the additional account and the exclusive
close cursor, so this is a distinct financial profile. Its validation replays
the retained calculation and every tax payment, then compares the complete
state. Changed balances, dates, loss results, payment histories and proposal
flags are rejected. Repeating the same accrual request produces the same pure
proposal; it does not establish globally unique durable posting. Later ordinary
return or lease events, financial-year carry, actual accounting admission and
annual completeness remain separate requirements.

\subsection{Continue ordinary events on the tax-adjusted book}

The additional tax payable must survive when ordinary business resumes. A
later return refund or lease payment has to use the cash left after tax
payments, and its journal must leave the tax balances intact unless that
method explicitly changes them. The continuation profile therefore retains
all 26 accounts, the return and lease memories, the assessment result and one
ordered event history. It extends the typed state and memory schemas directly;
it does not remove the tax payable to fit the older book.

The selected return methods instantiate the existing monetary definitions over
the complete account registry. Their memory operations still settle the return,
create the recovered lot or dispose of that lot, while preserving the additional
balances. The selected lease methods retain their original component equations
and map the resulting journal into the same registry. Tax expense and deferred
balances are unchanged by these ordinary methods. A separate, explicitly
instructed payment reduces income-tax payable and cash.

Assessment coverage constrains the next event. An ordinary event must occur on
or after the exclusive assessment endpoint and before the original financial
period ends. Otherwise it would change activity already used in the retained
tax calculation. A tax payment may occur on the accrual day, but cannot precede
an already completed operation. If assessment reaches the end of the declared
financial period, no ordinary continuation is possible without a separately
specified transition to the next opening.

\begin{samepage}
An earlier assessment provides a concrete same-period test. The author declares
the original two-year financial window, assesses the first interval through
31 December 2024 and places the later ordinary events on 1 January 2025.
The taxable expression of 2 and tax of $1/2$ are explicit hypothetical
assumptions. Paying that tax, settling the return, selling the recovered lot
and paying the lease use one cash balance:
\[
 500-\tfrac12-48+30-121=\tfrac{721}{2}.
\]
\end{samepage}
After the accompanying depreciation, recorded profit is $9/2$, assets are
$1401/2$, liabilities are 104 and equity is $1193/2$. Operating cash is
$-59/2$ and financing cash is $-110$. Direct cash, the indirect bridge,
the balance sheet and the equity movement reconcile exactly. The return's
later tax adjustment of $-24$ remains dated and identified for a subsequent
assessment; this continuation does not silently fold it into the earlier tax.

Payment order becomes consequential when cash is scarce. With opening cash
$485/4$, the lease payment of 121 alone leaves $1/4$. Paying tax of $1/2$
first leaves $483/4$, which cannot fund that lease payment. Reversing the order
allows the lease payment but leaves too little for tax settlement. Each request
stops at its first failed operation and retains the complete successful prefix.
It neither restores spent cash nor supplies financing automatically.

The result preserves the complete original report, loss state and conditional
tax calculation alongside the subsequent journals and gross cash movements.
Applying it recomputes the ordered prefix and compares the whole result.
Unknown findings remain unresolved, and repeated identities, changed balances
and backdated activity are refused. These checks establish selected same-period
execution under the declared assumptions. Further tax assessment, financial-year
carry and actual accounting admission remain separate work.

\subsection{Replace a provisional assessment without charging twice}

Later activity changes the question that the tax calculation must answer.
The continued book includes the return adjustment of $-24$, while its earlier
assessment covers only the preceding interval. A transition to the next
financial opening cannot make that unassessed activity disappear. The next
finite method therefore supplies an explicitly selected cumulative replacement:
a later calculation over the same original opening supersedes the earlier
provisional assessment. This is a hypothetical tax model. Independent tax
periods require a different method.

The replacement reads the actual continued 26-account book. It preserves the
cash already paid, retains the separate payable and includes every pretax
account and relevant lease right in the new declarations. Each native return
adjustment must appear exactly once with its event identity and amount.
Explicit additional assumptions must reconcile to the declared account totals,
and right-level tax bases must agree with their aggregate balances. Missing
declarations remain unresolved; incompatible scope, amounts or dates are refused.

The loss calculation returns to the original opening inventory. Using the
provisional closing ledger as a new opening could count a loss twice or make
it available inside the cumulative interval that generated it. The replacement
instead applies the declared usage, cap, eligibility and expiry rule once to
the original lots. The earlier result remains attached as superseded evidence.
For example, a provisional generated loss of 19 replaced by a cumulative loss
of 7 leaves 7, not 26. An original opening loss of 19, provisionally used by 1
and later used by 19 cumulatively, has total usage 19 rather than 20.

\begin{samepage}
Let $C_0$ be current-tax expense already recorded, $C_1$ the replacement total
after explicit loss usage, and $S$ actual tax payments. The additional expense
and resulting unpaid liability are
\[
 \Delta C=C_1-C_0,\qquad L_1=C_1-S.
\]
\end{samepage}
The method debits current-tax expense and credits tax payable by $\Delta C$.
A negative amount reverses those balances; it does not change cash. The
selected profile requires $L_1\geq0$. If earlier payments exceed the new total,
it preserves the calculation but refuses the financial proposal because a
tax receivable or refund needs a separately specified method.

In the paid mixed example, expense and payments are both initially $1/2$.
A declared replacement total of $21/4$ therefore adds only $19/4$ to expense
and payable. Cash stays at $721/2$. Pretax profit of 5 less total current-tax
expense of $21/4$ gives a recorded loss of $1/4$; assets of $1401/2$ less
liabilities of $435/4$ equal equity of $2367/4$. The operating cash result
remains $-59/2$. These amounts follow the explicit cumulative assumptions.

If no tax has been paid, replacing expense of $1/2$ with a total of $1/4$
reverses expense and payable by $1/4$ without moving cash. If $1/2$ was already
paid, the same replacement would leave a receivable of $1/4$ and is refused by
this profile. Both outcomes retain the complete continued financial history
and the replacement calculation for inspection.

The new coverage endpoint is later than the provisional endpoint and no later
than the original financial period boundary. The recorded events must fit
within it. Validation recomputes the continued book, the cumulative calculation,
the replacement losses and the incremental journal; applying a proposal compares
the whole result. Deferred-tax balances remain unchanged. Further assessment
chains, deferred-tax recognition, financial-period carry, source admission and
annual completeness still require their own methods and evidence.


\paragraph{Begin another period with recorded balances and unresolved findings.}
A conditional simulation can need another period before all the preceding
period's accounting questions have been resolved. For example, the scenario
above has cash, an unpaid current-tax liability, lease rights and a recorded
profit, but its deferred-tax assessment and the completeness of its historical
inputs remain open. Those recorded quantities are sufficient to construct a
provisional opening for selected subsequent transactions. They are insufficient
to declare that the preceding annual financial statements are complete.

The explicit model \texttt{provisional\_recorded\_balance\_carry\_v1} makes
that distinction part of execution. The request identifies the complete prior
reassessment, a unique carry event and the next period's start and exclusive
end. The prior current-tax assessment and loss history must reach the preceding
financial endpoint, and the next period must start at that endpoint. The prior
closing date must be its last included day. No calendar year or assessment of
missing transactions is inferred from those dates. Separate findings describe
event coverage, deferred-tax assessment, source admission and accounting
admission. Unknown and negative findings remain visible in the resulting state;
affirmative caller assumptions also leave annual readiness unestablished.

The transfer follows the book's debit-positive convention. Let $P$ denote the
recorded profit, $R$ the signed retained-earnings balance, and $b_k$ each
signed account balance immediately before the boundary. Let $\mathcal I$
contain the temporary income and expense accounts. The new opening satisfies
\[
 P=-\sum_{k\in\mathcal I}b_k,\qquad
 R^{\mathrm{open}}=R-P,\qquad
 b_k^{\mathrm{open}}=0\quad(k\in\mathcal I).
\]
Every other account keeps its recorded amount. A profit therefore increases the
credit balance in retained earnings; a loss reduces it. The transfer journal
sums to zero without a balancing plug. In the mixed scenario, recorded profit
is $-1/4$ and retained earnings is $-12$, so the opening signed balance is
$-12-(-1/4)=-47/4$. Cash stays $721/2$ and the unpaid tax liability stays
$19/4$. Positive-profit and zero-profit cases separately check the same sign
convention.

The boundary changes period accumulators without discarding economic memory.
Return components, recovered lots including disposed identities, lease rights,
measurement dates, loss quantities and lifetime event identifiers survive
exactly. The prior result retains its cash flows and tax adjustments together
with hashes of the closing book and report. The next period begins with empty
cash-flow and tax-adjustment lists, so an old receipt or deduction cannot be
counted again. Opening deferred-tax allocations copy the recorded asset and
liability; that copying does not assess their recognition or measurement. The
new period's assessment cursor denotes an empty interval and its current tax is
explicitly unassessed.

Subsequent statements use the actual carried opening. Paying the opening tax
liability of $19/4$ reduces cash from $721/2$ to $1423/4$, eliminates that
liability and leaves current-period tax expense at zero. In the indirect cash
calculation, the reduction in the opening payable explains the operating cash
outflow. It must not be charged to profit again. If depreciation of 10 follows,
recorded profit is $-10$, assets are $2743/4$, liabilities are 104 and equity is
$2327/4$. Depreciation is added back in the indirect bridge, leaving operating
cash flow at $-19/4$. The balance sheet, equity movement and direct and indirect
cash calculations reconcile to the same opening and closing state.

A second example carries an unsettled return component instead. Settlement
pays 48, releases a refund liability of 40 and a recovery asset of 16, and
recognizes recovered inventory of 20. Its cost-of-sales credit of 4 and selling
expense of 8 reduce profit by 4. Selling that inventory for 30 then earns 10,
so the two transactions together earn 6 and reduce cash by 18. Starting from
cash of 500, the following-period cash is 482. The indirect calculation gives
$6+16-40=-18$. The new return tax adjustment of $-24$ is retained for a later
assessment; it is not absorbed into an asserted tax-complete closing balance.

Ordered execution retains the completed prefix if a later event fails. A
payment still needs both an unpaid liability and actual shared cash; a reused
identity or an event outside the next period is refused. Applying a result
recomputes the entire predecessor and subsequent prefix and compares the
complete result, including unresolved findings and archived history. The
checked implementation covers this provisional boundary and selected return,
lease and tax-payment operations. New-period tax assessment, repeated carry,
complete event coverage and deferred-tax recognition require their own methods
and evidence before the scenario can support a complete annual close.


\paragraph{Measure deferred balances from explicit recognition assumptions.}
Carrying a recorded deferred-tax asset and liability preserves the book, but a
later assessment must explain whether their amounts should change. The mixed
book records an asset of 10 and a liability of 4. Its current-tax declarations
also contain tax bases. Neither fact establishes the amount of deductions that
can be recognized, the applicable rate on recovery or settlement, or whether a
movement belongs in profit or loss. The conditional deferred-assessment method
therefore requires those decisions separately and binds them to the actual
reassessed book.

The distinction follows the historical source treatment. IAS 12 paragraphs 7--9
distinguish asset and liability tax bases; paragraphs 24 and 27--29 condition
asset recognition on available taxable profit and its relevant grouping and
timing. Paragraphs 47 and 51 connect measurement to the applicable rate and
manner of recovery or settlement. Paragraphs 58 and 61A distinguish profit-or-loss
movements from tax related to items outside profit or loss
\citep[IAS 12, paras. 7--9, 24, 27--29, 47, 51, 58 and 61A]{eu20231803}.
These retained historical provisions explain the method's boundaries. The
example's rates, bases and recovery capacity remain explicit hypothetical
premises, rather than findings about an actual company's tax treatment.

The selected model uses one declared entity, tax authority and recovery group,
and keeps gross deferred assets and liabilities separate. Every nondeferred
asset and liability account is represented. Inventory is divided between other
inventory and retained lots; lease assets and liabilities are divided by right;
each remaining tax-loss lot is included separately. A component and its aggregate
cannot both contribute the same difference. Item bases must reconcile with the
complete prior declarations, while recovered-lot bases must also agree with
component memory. A finding about the absence of unrepresented tax-only items
and credits makes this finite scope explicit.

Let $C_i$ be the positive magnitude of an item's carrying amount, $B_i$ its
declared tax base and $r_i$ its declared nonnegative rate. An asset has signed
taxable difference $C_i-B_i$; a liability has $B_i-C_i$. A positive difference
is taxable and a negative one is deductible. For a loss lot, the supplied model
uses the negative of its remaining quantity, without consuming that quantity.
Write $d_i$ for the nonnegative deductible difference and $t_i$ for the
nonnegative taxable difference. The caller supplies a recognized deduction
$u_i$ and a common recovery capacity $K$, all in the same tax-base currency units:
\[
 0\leq u_i\leq d_i,\qquad \sum_i u_i\leq K,\qquad
 A_1=\sum_i r_i u_i,\qquad L_1=\sum_i r_i t_i.
\]
The sums include deductible temporary differences and loss lots together, so
one capacity cannot be reused for each item. The unrecognized amount $d_i-u_i$
remains visible. A deferred-tax liability does not automatically enlarge $K$:
the timing, entity and income restrictions needed to support recovery must be
supplied explicitly. The method does not produce that support by forecasting
future taxable profit.

Before changing balances, the request allocates the recorded asset $A_0$ and
liability $L_0$ completely across the same item population. Those allocations
must sum exactly to the book's recorded gross amounts. They explain inherited
balances rather than inventing a new assessment of the original opening.
For the supported profit-or-loss route, the change in deferred expense is
$(L_1-L_0)-(A_1-A_0)$. The journal debits the asset by $A_1-A_0$, credits the
liability by $L_1-L_0$, and records the balancing expense or benefit. Negative
changes reverse the corresponding balance. The request must support that
allocation; an equity, other-comprehensive-income or business-combination route
is refused by this particular implementation instead of being charged to profit.

In the checked example, property, plant and equipment has carrying amount 100
and tax base 60. At rate $1/4$, its taxable difference of 40 gives a liability
of 10. Debt has carrying amount 100 and tax base 60, so its difference is
deductible. Recognizing 32 of the available 40 gives an asset of 8 and leaves 8
unrecognized. Relative to the recorded asset of 10 and liability of 4, the asset
falls by 2 and the liability rises by 6; deferred expense is therefore 8.
Cash stays $721/2$, current tax and its unpaid liability are unchanged, and
recorded profit becomes $-33/4$. Assets of $1397/2$ equal liabilities of $459/4$
plus equity of $2335/4$. Adding the deferred expense back in the indirect cash
calculation leaves operating cash at $-59/2$.

Loss recognition uses the same capacity discipline. Recognizing deductions of
32 for the debt and 10 for a loss lot needs capacity of at least 42; capacity
of 32 is insufficient even if each item separately fits it. With capacity 42
and rate $1/4$, their combined deferred asset is $21/2$. The loss ledger still
contains the original remaining quantity of 19: recognizing a tax benefit does
not use the loss in the current-tax calculation. Matching bases, zero rates,
partial recognition and opposite asset/liability differences have separate
checked cases.

The proposal retains the entire prior result and changes only deferred balances,
deferred expense and its new event identity. Cash flows, current tax, loss
history, component state and assessment dates remain exact. Unknown necessary
findings stay unresolved, and an inconsistent population, allocation, base or
capacity produces no changed financial state. Pure apply compares the complete
recomputed proposal. Actual source and accounting admission, broader recovery
and allocation methods and annual completeness remain open. The preceding
frozen carry interface accepts a current-tax reassessment request; consuming
this newer result requires an explicit successor interface that preserves the
deferred journal and recognition schedule.


\paragraph{Carry the assessed book without renewing its recognition assumptions.}
The deferred journal changes the financial opening that a later period must
inherit. Replaying the older current-tax-only request would restore the asset
of 10 and liability of 4 and omit the deferred expense of 8. The resulting
retained earnings would also be wrong for the assessed scenario. The successor
carry therefore accepts the complete deferred-assessment request and recomputes
that proposal before transferring its actual closing balances.

The transfer uses the same debit-positive convention as the preceding carry.
Recorded profit is now $-33/4$, so retained earnings moves from $-12$ to
$-12-(-33/4)=-15/4$. The deferred asset of 8, deferred liability of 10, cash of
$721/2$ and current-tax payable of $19/4$ remain unchanged. Temporary income
accounts, including the deferred expense, begin at zero. Thus the next period
inherits the assessed economic amounts without charging the old expense again.

The complete predecessor remains attached. Its item-level differences,
recognized and unrecognized deductions, recorded allocations, support findings
and recovery capacity are retained with the original loss history. This archive
explains the opening balances; it does not renew the capacity or establish
new-period recognition. Opening deferred allocations copy the recorded 8 and
10, while both current and deferred tax in the next period are explicitly
unassessed. Unknown or negative carry findings remain visible independently of
the predecessor's conditional calculation.

If the next period pays the carried tax liability and records depreciation of
10, cash becomes $1423/4$ and recorded profit is $-10$. Assets of $2735/4$
equal liabilities of 110 plus equity of $2295/4$. Neither current nor deferred
tax expense is repeated. Operating cash is $-19/4$: depreciation is added back,
and the reduction in the opening tax payable explains the payment. The direct
cash calculation and the equity movement use this same carried opening.

The request retains the prior assessment and loss endpoints, requires a
contiguous next period and keeps lifetime event identities. Selected return,
lease and payment operations then act on the full new state. A failed operation
preserves the completed prefix. Applying a result compares the full recomputed
predecessor, opening, recognition archive and subsequent operations. The checked
scope remains provisional: a new current-tax assessment must distinguish
payments of the opening liability from payments of the new period's tax, and
new deferred recognition must reassess its own information. Repeated carry and
complete annual source and event coverage remain separate requirements.


\paragraph{Keep carried tax payments separate from the new assessment.}
The first current-tax assessment after carry has a different opening from the
same-period reassessment above. Current-period expense begins at zero, but the
book can still owe tax from the preceding period. Paying that opening liability
settles an old assessment; it does not reduce the new period's tax expense.
Calculating paid tax as current expense less the outstanding payable would
therefore be wrong for this opening.

The request binds the complete assessed-carry execution and begins its tax
interval at the new financial opening and loss cursor. It requires a completed
incoming execution and no previous assessment in the new period. Before that
first assessment, the selected payment kernel can consume only the carried
liability. Each payment can therefore be attributed to the prior assessment
from its exact event identity, date, amount and cash-flow record. A stopped
incoming execution is retained for inspection but is not silently assessed.

Let $P_0$ denote the opening tax payable, $S$ the completed payments against
that carried liability and $C$ the new conditional assessment. The two liability
records reconcile to the financial book as
\[
 P_{\mathrm{carried}}=P_0-S,\qquad
 P_{\mathrm{new}}=C,\qquad
 P_{\mathrm{total}}=P_0-S+C,\qquad 0\leq S\leq P_0.
\]
The first record retains the prior assessment's interval and identity with its
new-period settlements. The second records the new interval, assessment and
initially empty settlement history. The journal debits current-tax expense by
$C$ and credits the payable by the same amount. Cash, deferred balances and the
old settlements remain unchanged. Even a fully paid old liability keeps its
history rather than disappearing from the explanation.

For opening payable $19/4$ and new tax of 5, paying all $19/4$ leaves a total
payable of 5. Paying only 1 leaves $35/4$, and paying nothing leaves $39/4$.
The new expense is 5 in all three cases. In the fully paid example, after the
new-period depreciation of 10, recorded profit is $-15$, cash is $1423/4$,
assets are $2735/4$, liabilities are 115 and equity is $2275/4$. The increase
in the new payable cancels its expense in the indirect cash calculation; the
old payment remains the operating cash outflow of $19/4$.

Taxable profit is calculated from the new period's pretax accounts and explicit
adjustments. Each native return adjustment must appear exactly once, and the
complete tax-base declarations must reconcile at account and lease-right level.
Those declarations do not themselves assess deferred tax. In a return example,
new-period recorded profit of 6 and its native adjustment of $-24$ produce
taxable loss of 18. The method retains a new loss quantity of 18, available at
the interval's exclusive endpoint, without a refund or an automatic deferred
asset.

Loss use starts from the actual carried inventory. If a loss of 19 remains
available, explicit use against taxable profit of 20 leaves taxable profit of
1 and tax of $1/4$ at the declared rate. The history continues from the prior
cursor, with total use of 19. The selected close-day rule separately checks
expiry: a loss expiring before the usage day cannot be used, while a use before
expiry consumes only the declared quantity. Neither the previous deferred-tax
recognition nor its recovery capacity consumes or replenishes the loss inventory.

Applying the proposal compares the full predecessor, tax calculations, loss
movements, both liability records and financial statements. The checked scope
is a first assessment without credits, refunds or prepayments. Once both tax
vintages exist, subsequent payments need explicit allocation; the original
aggregate payment interface no longer supplies enough information. New-period
deferred assessment, repeated current assessments and later carry require their
own compatible methods. Source admission and complete annual event coverage
remain unresolved.


\subsubsection*{Name the tax liability settled by each payment}
Once carried tax and newly assessed tax are both payable, the payment amount
alone cannot identify which liability has been settled. A payment of 3 can reduce
either the carried payable of $19/4$ to $7/4$, leaving the new payable of 5
unchanged, or reduce the new payable to 2, leaving the carried payable at $19/4$.
Both choices reduce cash and total payable by 3. Their remaining tax obligations
have different histories, which the simulator must preserve.

The continuation therefore requires the caller to name each tax vintage and its
allocated amount. Every amount must be positive, no greater than the named
outstanding balance, and the allocations must sum to the payment. Actual shared
cash must cover that sum. An unknown authorization, a duplicated vintage, an
unknown liability or an excessive allocation prevents that operation. The
method supplies no first-in-first-out convention or proportional split.

For example, take the hypothetical first assessment through 30 June 2026 inside
the financial period ending on 31 December. Carried tax of $19/4$ is still
unpaid, new tax is 5, cash is $721/2$, and recorded profit is $-15$. Allocate 1
of a payment to carried tax and 2 to new tax. Cash becomes $715/2$; the two
payables become $15/4$ and 3, totaling $27/4$. Tax expense remains 5 and profit
remains $-15$. The payment is recorded in both affected settlement histories,
with their distinct amounts and the same dated cash movement.

An ordinary event can follow only from 1 July, the assessment's exclusive
endpoint. A further right-of-use depreciation charge of 2 reduces its asset from
80 to 78 and profit to $-17$, without changing either tax payable. Assets then
total $1367/2$, liabilities $467/4$, and equity including profit $2267/4$.
Operating cash flow remains $-3$. This depreciation is explicitly listed as
activity after the assessment; it has not received a new tax assessment.

The full book, prior loss inventory, deferred balances and unresolved findings
remain attached to the result. If a later operation fails, earlier completed
events and their settlement histories survive as a recorded prefix. Payments
can occur on the recorded assessment day, but an assessment at the financial
period's endpoint leaves no later date for ordinary events in that period.
Cumulative reassessment, new-period deferred recognition and another period
carry require their own compatible methods. The present calculation supplies
conditional bookkeeping evidence rather than annual accounting acceptance.


\subsubsection*{Replace the assessment while retaining the payments}
After further company activity, the new period's cumulative tax calculation may
change. Its recorded payments do not change with it. A replacement must therefore
start from the actual continued book and distinguish what was paid against the
carried liability from what was paid against the current period's assessment.
Using current expense less all tax payments would wrongly treat a payment of
last year's tax as a reduction in this year's expense.

Let $C_0$ be recorded current-period tax expense, $C_1$ the replacement cumulative
amount, $S$ the actual payments allocated to that current vintage, and $O$ the
remaining carried liability. The replacement records only the expense difference:
\[
  \Delta C=C_1-C_0,\qquad
  P_{\mathrm{current}}=C_1-S,\qquad
  P_{\mathrm{total}}=O+C_1-S.
\]
The journal debits current-tax expense by $\Delta C$ and credits payable by the
same signed amount. Cash and the carried liability remain unchanged. The current
vintage retains the identity referenced by earlier payments; its old assessment
and settlement record remain available alongside the replacement's provenance.

Continue the split-payment example above. Recorded current expense is 5,
current-vintage payments are 2, and carried payable is $15/4$. If the replacement
cumulative tax is 7, the expense increment is 2, current payable becomes 5 and
total payable becomes $35/4$. Cash stays at $715/2$ and recorded profit falls
from $-17$ to $-19$. Assets remain $1367/2$, while liabilities become $475/4$
and equity including profit becomes $2259/4$. If cumulative tax instead falls
to 3, the expense reduction is 2 and total payable becomes $19/4$.

A positive aggregate liability can conceal a different representation problem.
If cumulative current tax falls to 1 after payments of 2, the current vintage
would have a credit of 1 even though the carried liability leaves a positive
aggregate of $11/4$. The present positive-payable model refuses that replacement
and preserves the preceding book. A credit or refund needs a specified account
and treatment; it cannot be silently netted against the older obligation.

Losses are also recalculated cumulatively. Suppose the carried opening contains
an available loss of 19 and the first assessment uses all 19 against taxable
profit of 20. At a declared rate of one quarter, current tax is $1/4$. A
replacement using the same cumulative profit and deduction again yields $1/4$:
it restarts from the original opening inventory, rather than using the provisional
closing balance of zero or recording a second consumption of 19. Previous
provisional calculations remain in the history. If that loss expires before the
new assessment's usage date, the replacement must apply the stated expiry rule;
it cannot retain the earlier deduction automatically.

The cumulative interval must cover every completed event and cannot shorten the
previous assessment. A correction at the same endpoint is possible when no later
event lies outside it. Native return adjustments enter the replacement exactly
once; a worked return sequence has pretax income 6 and an adjustment of $-24$,
producing a new loss of 18 without an automatic deferred asset. Fresh deferred
recognition, another continuation or repeated replacement requires a compatible
method. Source admission, complete fiscal coverage and annual acceptance remain
unresolved.
